\section{Introduction}

Statistic is the foundation for all data-driven decision-making processes, also
in the AI era: in fact, nowdays machine learning models rely mostly on statistical
methods. The main difference resides in the purpose: \textit{ML models} are designed to 
make the most accurate prediction possible, \textit{statistical models} instead aim to 
make inference about the relationship between the data.
\begin{figure}[h]
    \centering
    \includegraphics[width=0.3\textwidth]{introduction/ml_statistic_intersection.png}
\end{figure}

In this course we are going to deal with \textit{multivariated data}, meaning that each 
datapoint contains several variables, defining its \textit{dimensions}. Representing this data 
with matrices is the more convient way.
\begin{figure}[h]
    \centering
    \includegraphics[width=0.3\textwidth]{introduction/multivariated_data_example.png}
\end{figure}

\vspace{1cm}

As for ML, we have different classes of methods based on the available data and how we 
interpret them:
\begin{description}[itemsep=0.6cm]
    \item[Supervised learning] \hfill \\
    \begin{minipage}{0.9\textwidth}
        \begin{wrapfigure}{l}{0.3\textwidth}
            \includegraphics[width=0.25\textwidth]{introduction/target_feature_relationship.png}
        \end{wrapfigure}
        \hfill \\
        In this class we can identify a \textbf{target} variable $Y$ and a set of "input"
        variables $X$ called \textbf{features}.  Targets and features relationship is
        expressed trougth a function affected by a certain error $\varepsilon$: 
        $Y=f(X)+\varepsilon$. The error $\varepsilon$ is something we can never get rid of
        it, we should think about it as a property of the data.\\
        Typically in this set of method we try to make predictions about new target values or 
        to find those features that are more important to explain the targets.
    \end{minipage}
    
    \item[Unsupervised learning] \hfill \\
    \begin{minipage}{0.9\textwidth}
        \begin{wrapfigure}{l}{0.3\textwidth}
            \includegraphics[width=0.25\textwidth]{introduction/example_unsupervised.png}
        \end{wrapfigure}
        \hfill \\
        In this case there aren't target variables, but all the variables are features 
        ($f(X)\approx X$). \\ This set of models focuses on exploiting structures and patterns
        within the data itself.
    \end{minipage}
\end{description}
\vspace{0.8cm}

\begin{figure}[t]
    \centering
    \includegraphics[width=0.6\textwidth]{introduction/methods_taxonomy.png}
    \caption{Taxionomy of the methods explored during the course.}
\end{figure}
\newpage

\subsection{Model goodness}
How do we find a good model for our data? \\
To start, let's take linear models as examples, like $Y=\beta_0 x_0 + \beta_1 x_1 \dots + 
\beta_p x_p$. \\
The models we want to use are expressed through a set of \textbf{parameters} of our choice
(in the previous example we have $p+1$ parameters); the goal is to tune these parameters
using training data in order to shape our function accordingly to the data. However most 
of the times linear models are not complex enough, and we have to shift toward more complex 
ones.
\begin{figure}[h]
    \centering
    \includegraphics[width=0.65\textwidth]{introduction/linear_quadratic_models_comparison.png}
    \caption{Comparison between linear model ($f(X)=\beta_o + \beta_1 x$) and quadratic 
    model ($f(X)= \beta_0 + \beta_1 x +\beta_2 x^2$).}
\end{figure}

Usually the choice of a model is influenced by some trade-offs:
\begin{itemize}
    \item \textit{prediction accuracy vs interpretability} : we might trade-off the model's
    accuracy to make it more interpretable
    \item \textit{parsimony principle} : usually we prefer simpler models involving fewer
    variables since they're easier to manage
    \item \textit{good fit vs over-fit or under-fit } : see below
\end{itemize}
\newpage

\subsection{Underfitting vs Overfitting problem}

Start with the example of estimating the height of the next student entering the room 
during lecture, in order to do that our prediction $\hat{y}$ should minimize a certain 
function, usually the \textbf{Mean Squared prediction Error (MSE)}.
\begin{equation*}
    MSE(\hat{y})=\frac{1}{n}\sum_{i=1}^n (\hat{y}_i-y_i)^2
\end{equation*}

If we know the height of all students, we want to use the population mean $\mu$ as 
predictor, doing so we get as MSE $\mathbb{E}[(Y-\mu)^2]=Var(Y)=\sigma^2$.\\
\begin{center} \textcolor{red}{
    \textbf{The data variance is irreducible, we can't do better than that.}   
} \end{center}
\hfill \\

However, it's impossible to have information about the whole population; so what we do
is taking a small sample from it over and over, using the sample mean $\bar{Y}$ as $\mu$
and getting $MSE_{avg}=\mathbb{E}[(\hat{Y}-Y)^2]$. By decomposing the average MSE we are
able to understand which are our sources of error: this reduction leads to two source of
errors, \textit{bias} and \textit{variance}.
\begin{description}
    \item[Bias]: $Bias(\hat{Y})=\expect[\hat{Y}-\mu]$, we have bias if our predictions 
    are systematically to high or too low
    \item[Variance]: $Var(\hat{Y})=\expect[(\hat{Y}-\expect[\hat{Y}])^2]$ , we have
    high variance if our predictions are very sensitive on the sampling we draw from 
    the universe
\end{description}
\begin{figure}[h]
    \centering
    \includegraphics[width=0.6\textwidth]{introduction/bias_variance_idea.png}
\end{figure}

So if we try to fit a model $f(X)$ with some training data, and we end up with $\hat{f}(X)$
with some test observation $(x_0,y_0)$, we have that
\begin{equation*}
    \expect[y_0-\hat{f}(x_0)] = Var[\hat{f}(x_0)] + Bias[\hat{f}(x_0)]^2 + Var(\varepsilon)
\end{equation*}
\begin{description}
    \item[Variance]: captures how much our model changes based on the training dataset
    \item[Bias]: the inherent error we get from our models
    \item[Noise]: ambiguity due the data distribution   
\end{description}
\newpage

If we take a model too much simple, we end up with something having high bias, but not really
susceptible to training set (low variance) $\Rightarrow$ \textbf{Underfitting}.
\begin{figure}[h]
    \centering
    \includegraphics[width=0.6\textwidth]{introduction/underfitting.png}
\end{figure}

If instead we use a model too complex we get something with very low bias, but extrimelly
sensitive to the data we sample $\Rightarrow$ \textbf{Overfitting}.
\begin{figure}[h]
    \centering
    \includegraphics[width=0.6\textwidth]{introduction/overfitting.png}
\end{figure}

What we see in practise is that Bias and Variance are strictly correlated: while one of the
two decreases, the other increases. This force us to find some trade-off between these twos.
\begin{figure}[h]
    \centering
    \includegraphics[width=0.6\textwidth]{introduction/tradeoff_bias_variance.png}
\end{figure}
